\section{Dense-gap calculations}
\label{app:dense-gap-calculations}

This appendix supplies the parameter choices and gap calculations
for the proof sketches in the main text. All parameters are fixed
once the excluded approximation factor is fixed.

\subsection{UGC treedepth}
\label{app:dense-ugc-td}

\begin{proof}[Dense treedepth gap]
Use the coefficients $A_{q,\eta}$ and $B_{q,\eta}$ from
\eqref{eq:ugc-yes} and~\eqref{eq:ugc-no}, with $n=|V(G)|$ for
the independent-set source graph $G$.
At zero source error, the ratio of the NO and YES coefficients is
$B_{q,0}/A_{q,0}=4q/(3(q+1))$, which approaches $4/3$ as $q$
increases. Fix $1\leq r<4/3$ and choose
$q=\lfloor3r/(4-3r)\rfloor+1$. Put
$\delta=[q(4-3r)-3r]/6>0$ and
$C=2q/3+r(q-1)>0$. Choose
$0<\eta<\min\{1/10,\delta/(2C)\}$. Then
\begin{equation}\label{eq:ugc-gap-parameter}
 \left(\frac{2q}{3}+r(q-1)\right)\eta
 <\frac{q(4-3r)-3r}{6}.
\end{equation}
Then
\[
B_{q,\eta}-rA_{q,\eta}
 =\frac{2q}{3}(1-\eta)
   -r\left(\frac{q+1}{2}+(q-1)\eta\right)
  =\frac{q(4-3r)-3r}{6}
   -\left(\frac{2q}{3}+r(q-1)\right)\eta
   \overset{\eqref{eq:ugc-gap-parameter}}{>}0.
\]
An $r$-approximation returning a feasible forest has height at most
$rA_{q,\eta}n$ in a YES instance. In a NO instance every feasible
forest has height greater than $B_{q,\eta}n>rA_{q,\eta}n$ by
\eqref{eq:ugc-no} and~\eqref{eq:ugc-gap-parameter}. Comparing the
returned height with $rA_{q,\eta}n$ therefore distinguishes the cases.
Since $q$ and $\eta$ depend only
on $r$, this is a polynomial reduction.
\end{proof}

\subsection{UGC total elimination depth}
\label{app:dense-ugc-total}

\begin{proof}[Dense total-depth gap]
Fix $1\leq r<8/7$. For an integer $q\geq2$ and a source error
$0<\eta<1/10$, let $G$ be the independent-set source graph and
put $n=|V(G)|$. The source thresholds are
\begin{equation}\label{eq:ugc-total-depth-source-thresholds}
\textup{YES:}\quad \alphaI(G)
 \overset{\text{Thm.~\ref{thm:bounded-degree-kr-source}}}{\geq}(1/2-\eta)n,
 \qquad \textup{NO:}\quad \vc(G)
 \overset{\eqref{eq:cover-independent-complement}}{=}n-\alphaI(G)
 \overset{\text{Thm.~\ref{thm:bounded-degree-kr-source}}}{>}(1-\eta)n.
\end{equation}
For $0\leq\eta<1/10$, define the coefficient polynomials
$A^{\Sigma}_{q,\eta}=(q^2-q(q-1)(1/2-\eta)^2)/2$ and
$B^{\Sigma}_{q,\eta}=(3/7)q^2(1-\eta)^2$.
The value $\eta=0$ is used only for the limiting calculation;
the reduction uses a positive source error.
The corresponding output satisfies
\[
\begin{aligned}
 \textup{YES:}\quad \totd(S_q(G))
 &\overset{\text{Lem.~\ref{lem:total-depth-completeness}}}{\leq}
 \frac{qn(qn+1)-q(q-1)\alphaI(G)^2}{2}
  \overset{\eqref{eq:ugc-total-depth-source-thresholds}}{\leq}
 A^{\Sigma}_{q,\eta}n^2+qn/2,\\
 \textup{NO:}\quad \totd(S_q(G))
 &\overset{\eqref{eq:total-depth-structural}}{\geq}
 \frac37q^2\vc(G)^2+\frac{qn}{2}
 \overset{\eqref{eq:ugc-total-depth-source-thresholds}}{\geq}
 B^{\Sigma}_{q,\eta}n^2+qn/2.
\end{aligned}
\]
At $\eta=0$, the ratio of the quadratic NO and YES coefficients
is $24q/(7(3q+1))$, which approaches $8/7$ as $q$ increases.
Choose a fixed integer $q$ satisfying
\begin{equation}\label{eq:ugc-total-depth-q-choice}
 q>\frac{7r}{24-21r}.
\end{equation}
At $\eta=0$,
\[
\begin{aligned}
 A^{\Sigma}_{q,0}
 &=\frac{q^2-q(q-1)/4}{2}=\frac{q(3q+1)}8,\\
 B^{\Sigma}_{q,0}-rA^{\Sigma}_{q,0}
 &=\frac{3q^2}{7}-\frac{rq(3q+1)}8
   =\frac{q\bigl(24q-7r(3q+1)\bigr)}{56}
   \overset{\eqref{eq:ugc-total-depth-q-choice}}{>}0.
\end{aligned}
\]
Indeed,
$24q-7r(3q+1)=(24-21r)q-7r
\overset{\eqref{eq:ugc-total-depth-q-choice}}{>}0$.
To choose the source error explicitly, put
\[
 D_0=\frac{q\bigl((24-21r)q-7r\bigr)}{56},\qquad
 K=\frac{6q^2}{7}+\frac{rq(q-1)}2,\qquad
 L=\frac{3q^2}{7}+\frac{rq(q-1)}2.
\]
Here $D_0>0$ by~\eqref{eq:ugc-total-depth-q-choice}, and $K,L>0$.
Choose a fixed
$0<\eta<\min\{1/10,D_0/(2K)\}$. Expanding the two coefficients gives
\begin{equation}\label{eq:ugc-total-depth-coefficient-gap}
B^{\Sigma}_{q,\eta}-rA^{\Sigma}_{q,\eta}
 =D_0-K\eta+L\eta^2
 \overset{K\eta<D_0/2}{>}\frac{D_0}{2}>0,\qquad B^{\Sigma}_{q,\eta}>rA^{\Sigma}_{q,\eta}.
\end{equation}
For a source graph $G_0$ on $n_0\geq1$ vertices, put
\[
 m=1+\left\lfloor
 \frac{q(r-1)}{2(B^{\Sigma}_{q,\eta}-rA^{\Sigma}_{q,\eta})}
 \right\rfloor,
 \qquad G=\mathop{\dot\bigcup}_{j=1}^{m}G_0^{(j)},
\]
where $G_0^{(1)},\ldots,G_0^{(m)}$ are disjoint copies of $G_0$.
Then the order $n=mn_0\geq m$ satisfies
\begin{equation}\label{eq:ugc-total-depth-source-size}
 n>\frac{q(r-1)}
 {2(B^{\Sigma}_{q,\eta}-rA^{\Sigma}_{q,\eta})}.
\end{equation}
The denominator is positive by~\eqref{eq:ugc-total-depth-coefficient-gap}.
Since there are no edges between the copies,
\[
 |V(G)|=m|V(G_0)|,\qquad
 \alphaI(G)=m\alphaI(G_0),\qquad
 \Delta(G)=\Delta(G_0).
\]
Thus the source promise remains unchanged. The YES and NO
bounds then have gap
\[
\left(B^{\Sigma}_{q,\eta}n^2+\frac{qn}{2}\right)
 -r\left(A^{\Sigma}_{q,\eta}n^2+\frac{qn}{2}\right)=\left(B^{\Sigma}_{q,\eta}-rA^{\Sigma}_{q,\eta}\right)n^2
   -\frac{(r-1)qn}{2}
   \overset{\eqref{eq:ugc-total-depth-source-size}}{>}0.
\]
The reduction has gap greater than $r$.
All parameters, including $m$, depend only on $r$, so the reduction
is polynomial.
\end{proof}

\subsection{UGC treewidth}
\label{app:dense-ugc-tw}

\begin{proof}[Dense treewidth gap]
As the source error decreases and $R$ increases, the leading NO
and YES coefficients approach $2/3$ and $1/2$, giving ratio $4/3$.
Fix $1\le r<4/3$ and put $\delta=2/3-r/2>0$. Choose
$0<\zeta<\min\{1/12,\delta/[2(1+r)]\}$ and then an integer
$R>2r/(3\delta)$. These choices give
$(1+r)\zeta+r/(3R)<\delta$, or equivalently
\begin{equation}\label{eq:ugc-tw-parameter-gap}
 2/3-\zeta>r(1/2+\zeta+1/(3R)).
\end{equation}
Here $\zeta$ controls the source error and $1/(3R)$ controls the
additional YES cost of the replication construction.
Take the $N$-vertex source graph $G$ from Theorem~\ref{thm:ugc-tripartite-cover}
with $k=3$. Its parts have size $N/3$. Since $\zeta<1/12$, both
cover thresholds, $(1/2+\zeta)N$ and $(2/3-\zeta)N$, exceed $N/3$.
Thus the $n_{\max}=N/3$ term in the maximum of
Lemma~\ref{lem:tw-replication} does not weaken either case, and the
lemma and source promise give
\[
\begin{aligned}
 \textup{YES:}\quad \operatorname{tw}(K_R(G))
 &\overset{\text{Lem.~\ref{lem:tw-replication}}}{\leq}
 R\max\{\vc(G),N/3\}+N/3-1
  \overset{\text{Thm.~\ref{thm:ugc-tripartite-cover}}}{\leq}
 R(1/2+\zeta)N+N/3-1,\\
 \textup{NO:}\quad \operatorname{tw}(K_R(G))
 &\overset{\text{Lem.~\ref{lem:tw-replication}}}{\geq}
 R\vc(G)-1
 \overset{\text{Thm.~\ref{thm:ugc-tripartite-cover}}}{\geq}
 R(2/3-\zeta)N-1.
\end{aligned}
\]
The choices of $\zeta,R$ make
$\gamma:=R(2/3-\zeta)-r\bigl(R(1/2+\zeta)+1/3\bigr)
\overset{\eqref{eq:ugc-tw-parameter-gap}}{>}0$.
The NO lower bound minus $r$ times the YES upper bound is
$\gamma N+r-1\overset{\eqref{eq:ugc-tw-parameter-gap}}{>}0$.
Adding one to both treewidth bounds gives the
maximum-bag-size bounds, whose corresponding difference is
$\gamma N\overset{\eqref{eq:ugc-tw-parameter-gap}}{>}0$.
Thus both objectives have gap greater than $r$.
More explicitly, let $U=R(1/2+\zeta)N+N/3-1$ be the displayed
YES width bound. A factor-$r$ approximation returning a tree
decomposition has width at most $rU$ in the YES case, whereas
every decomposition has width greater than $rU$ in the NO case.
For maximum bag size, use the cutoff $r(U+1)$ instead.
The construction is polynomial, and its three specified cliques
each have size $RN/3$.
\end{proof}

\subsection{SSE treedepth}
\label{app:dense-sse-td}

\begin{proof}[Dense treedepth gap]
Let $c>0$ be the absolute constant in
Lemma~\ref{lem:sse-unweighted-gap}.
Using the convenient YES coefficient $4\varepsilon$, the ratio
of the NO and YES coefficients is $c/(8\sqrt\varepsilon)$,
which grows without bound as $\varepsilon$ decreases.
Fix $C\geq1$ and choose a fixed $\varepsilon>0$ small enough
that Lemma~\ref{lem:sse-unweighted-gap} applies,
$c\sqrt\varepsilon\leq1/4$, and $\sqrt\varepsilon<c/(8C)$.
Since $\varepsilon$ is fixed, $q=\lceil2/\varepsilon\rceil$ is a
constant. Apply that lemma with this $q$ to obtain an
unweighted host $G^*$ with $m$ edges, and output $\mathcal L(G^*)$.

In the YES case, let $D$ be the multiset of host edges joining different
partition classes. The host bounds are
\begin{equation}\label{eq:sse-yes-host-bounds}
|D|
 \overset{\text{Lem.~\ref{lem:sse-yes-edge-counts}}}{\leq}
 m\max_{i\in[q]}\Phi_{G^*}(V_i)
 \overset{\text{Lem.~\ref{lem:sse-unweighted-gap}}}{\leq}3\varepsilon m,\qquad |E(G^*[V_i])|
 \overset{\text{Lem.~\ref{lem:sse-yes-edge-counts}}}{\leq}
 \vol_{G^*}(V_i)/2
 \overset{\text{Lem.~\ref{lem:sse-unweighted-gap}}}{\leq}(1+\varepsilon)m/q.
\end{equation}
Consequently
\begin{equation}\label{eq:sse-yes-height}
\begin{aligned}
 \td(\mathcal L(G^*))
 &\overset{\text{Lem.~\ref{lem:sse-complete}}}{\leq}
 |D|+\max_{i\in[q]}|E(G^*[V_i])|\\
 &\overset{\eqref{eq:sse-yes-host-bounds}}{\leq}
 3\varepsilon m+(1+\varepsilon)m/q
 \leq3\varepsilon m+(1+\varepsilon)\varepsilon m/2
 =(7/2+\varepsilon/2)\varepsilon m<4\varepsilon m,
\end{aligned}
\end{equation}
where $q=\lceil2/\varepsilon\rceil$ gives $1/q\leq\varepsilon/2$,
and $\varepsilon<1$ gives $7/2+\varepsilon/2<4$.
In the NO case, $(c/2)\sqrt{\varepsilon}\leq1/8<1/4$, so
Corollary~\ref{lem:sse-sound} gives
\[
 \td(\mathcal L(G^*))
 \overset{\text{Lem.~\ref{lem:sse-unweighted-gap}, Cor.~\ref{lem:sse-sound}}}{\geq}
 (c/2)\sqrt\varepsilon\,m
 >4C\varepsilon m,
\]
by the choice of $\varepsilon$. The NO bound therefore exceeds the
YES bound by a factor greater than $C$.
Indeed, a $C$-approximation returning an elimination forest has height
at most $4C\varepsilon m$ in the YES case, whereas every valid forest
has larger height in the NO case. Comparing the returned height
with this threshold distinguishes the source promise.
The threshold uses only the fixed parameters and the observable
output order $m$; all construction sizes are polynomial for these
fixed parameters.
\end{proof}

\subsection{NP-hardness of treedepth}
\label{app:dense-np-td}

\begin{proof}[Dense treedepth gap]
The leading NO and YES coefficients below approach $9$ and $17/2$
as the source error decreases, giving ratio $18/17$.
Fix $1\le r<18/17$ and put $\delta_0=9-17r/2>0$.
Choose a fixed
$0<\epsilon<\min\{1/4,\delta_0/(2(1+r))\}$.
Then $9-\epsilon-r(17/2+\epsilon)
=\delta_0-(1+r)\epsilon>\delta_0/2$.
Put $\gamma=9-\epsilon-r(17/2+\epsilon)>0$. Obtain a source system
$F$ from Theorem~\ref{thm:xor-source} with error $\epsilon$, and fix
$t=\lfloor(r-1)/\gamma\rfloor+1$. Replace $F$ by $t$ disjoint
copies on disjoint variable sets, and let $m$ denote their total
equation count. Since the original count is at least one,
\begin{equation}\label{eq:xor-td-replication-choice}
 \gamma m>r-1.
\end{equation}
Both the equation count and the optimum are multiplied by the fixed
integer $t$, which depends only on $r,\epsilon$. Thus replication
is polynomial and preserves the source promise and~\eqref{eq:xor-normalization}.
Let $H=H_{3m}(G_F)$ be the output of
Proposition~\ref{prop:xor-treedepth-identity}.
By Proposition~\ref{prop:xor-treedepth-identity} and the source
promise, the two cases satisfy
\[
\begin{aligned}
\textup{YES: }\td(H)
 &\overset{\eqref{eq:xor-treedepth-identity}}{=}
   19m/2-\OPT(F)+1
  \overset{\text{Thm.~\ref{thm:xor-source}}}{\le}
   19m/2-(1-\epsilon)m+1
   =(17/2+\epsilon)m+1,\\
\textup{NO: }\td(H)
 &\overset{\eqref{eq:xor-treedepth-identity}}{=}
   19m/2-\OPT(F)+1
  \overset{\text{Thm.~\ref{thm:xor-source}}}{\ge}
   19m/2-(1/2+\epsilon)m+1
   =(9-\epsilon)m+1.\end{aligned}
\]
The NO lower bound minus $r$ times the YES upper bound is
$\gamma m+1-r\overset{\eqref{eq:xor-td-replication-choice}}{>}0$,
proving the gap.
\end{proof}

\subsection{NP-hardness of treewidth}
\label{app:dense-np-tw}

\begin{proof}[Dense treewidth gap]
Fix $1\le r<12/11$ and put $\delta_0=6-11r/2>0$.
Choose a fixed
$0<\epsilon<\min\{1/4,\delta_0/(2(1+r))\}$, so that
$6-\epsilon-r(11/2+\epsilon)>\delta_0/2$. Put
$\delta=6-\epsilon-r(11/2+\epsilon)>0$, and set the fixed integer
$R=\lfloor3r/\delta\rfloor+1$, so that
\begin{equation}\label{eq:np-treewidth-replication-choice}
 R\delta>3r.
\end{equation}
These choices preserve the source separation and control the
additional bag size in the replicated decomposition.
Use Theorem~\ref{thm:xor-source} with error $\epsilon$ to obtain
a system $F$ with $m\ge1$ equations and occurrence bound $b_\epsilon$.
Apply the witness construction of
Section~\ref{subsec:np-treedepth}. By~\eqref{eq:xor-equal-parts} and
Lemma~\ref{lem:xor-cover}, the resulting simple unweighted graph
$G_F$ satisfies
\[
 |V(G_F)|=9m,\qquad |A|=|B|=|C|=3m,\qquad
 \Delta(G_F)\leq3b_\epsilon,
\]
and
$\vc(G_F)\overset{\text{Lem.~\ref{lem:xor-cover}}}{=}13m/2-\OPT(F)$.
The two graph transformations are
\[
 F\xrightarrow{\text{witness graph}}G_F
 \xrightarrow{\text{replicate and complete}}K_R(G_F).
\]
Since $\OPT(F)\leq m$, every instance satisfies
\[
 \vc(G_F)\overset{\text{Lem.~\ref{lem:xor-cover}}}{=}
 \frac{13m}{2}-\OPT(F)\geq\frac{13m}{2}-m=\frac{11m}{2}>3m.
\]
The source promise gives the two bounds
\begin{equation}\label{eq:np-treewidth-cover-gap}
\begin{aligned}
 \textup{YES:}\quad
 \vc(G_F)&\overset{\text{Lem.~\ref{lem:xor-cover}}}{=}
    \frac{13m}{2}-\OPT(F)
  \overset{\text{Thm.~\ref{thm:xor-source}}}{\leq}
    \frac{13m}{2}-(1-\epsilon)m
 =\left(\frac{11}{2}+\epsilon\right)m,\\
 \textup{NO:}\quad
 \vc(G_F)&\overset{\text{Lem.~\ref{lem:xor-cover}}}{=}
    \frac{13m}{2}-\OPT(F)
  \overset{\text{Thm.~\ref{thm:xor-source}}}{\geq}
    \frac{13m}{2}-\left(\frac12+\epsilon\right)m
 =(6-\epsilon)m.\end{aligned}
\end{equation}
Applying Lemma~\ref{lem:tw-replication} gives
\[
 \begin{aligned}
 \textup{YES:}\quad\operatorname{tw}(K_R(G_F))
   &\overset{\text{Lem.~\ref{lem:tw-replication}}}{\leq}
       R\vc(G_F)+3m-1
  \overset{\eqref{eq:np-treewidth-cover-gap}}{\leq}
       R(11/2+\epsilon)m+3m-1,\\
 \textup{NO:}\quad\operatorname{tw}(K_R(G_F))
   &\overset{\text{Lem.~\ref{lem:tw-replication}}}{\geq}
       R\vc(G_F)-1
 \overset{\eqref{eq:np-treewidth-cover-gap}}{\geq}
       R(6-\epsilon)m-1.
 \end{aligned}
\]
For every $m\ge1$, the NO bound minus $r$ times the YES bound is
\[
 (R\delta-3r)m+(r-1)
 \overset{\eqref{eq:np-treewidth-replication-choice}}{>}0,
\]
by the choice of $R$ and $r\ge1$. Thus the bounds are separated by
factor $r$ even with the additive constants $-1$.
The output has $9Rm$ vertices, and its three global cliques have
equal size $3Rm$. 
\end{proof}

\subsection{NP-hardness of total elimination depth}
\label{app:dense-np-total}

\begin{proof}
Fix $1\leq r<63/62=189/186$. Choose a fixed
\[
0<\epsilon<\min\left\{\frac14,
                  \frac{189-186r}{6(1+r)}\right\},\qquad \gamma:=(189-6\epsilon)-r(186+6\epsilon)>0.
\]
Apply Theorem~\ref{thm:xor-source} with this error, obtaining a system
with $m_0\geq1$ equations. Fix an integer $K>6(r-1)/\gamma$, and
take $K$ disjoint copies with fresh variable names. Denote the resulting
system by $F$; its total equation count $m=Km_0$ satisfies
\begin{equation}\label{eq:np-total-depth-source-scale}
 3(r-1)<\gamma m/2.
\end{equation}
This preserves its optimum fraction,
the occurrence bound, and~\eqref{eq:xor-normalization}.
Construct $G=G_F^\Sigma$ and put $N=99m$. Fix an integer
$C_r\geq9900r/\gamma$ and set
$M=C_rm$.
Both $K$ and $C_r$ depend only on $r,\epsilon$. Since $m\geq1$,
$N(N+1)=9801m^2+99m\leq9900m^2$, so these choices give
\begin{equation}\label{eq:np-total-depth-amplification-budget}
 N(N+1)\leq9900m^2,\qquad
 M\geq\frac{rN(N+1)}{\gamma m}.
\end{equation}
For the fixed parameters,
$M=\Theta(m)$. Proposition~\ref{prop:np-total-depth-parity-identity}
and Lemma~\ref{lem:np-total-depth-added-vertices} give
\begin{equation}\label{eq:np-total-depth-W-gap}
\begin{aligned}
 \textup{YES:}\quad W(G)
 &\overset{\eqref{eq:np-total-depth-parity-W}}{=}192m-6\OPT(F)+3
 \overset{\text{Thm.~\ref{thm:xor-source}}}{\leq}(186+6\epsilon)m+3,\\
 \textup{NO:}\quad W(G)
 &\overset{\eqref{eq:np-total-depth-parity-W}}{=}192m-6\OPT(F)+3
 \overset{\text{Thm.~\ref{thm:xor-source}}}{\geq}(189-6\epsilon)m+3.
\end{aligned}
\end{equation}
Consequently,
\[
\begin{aligned}
 \textup{YES:}\quad \totd(H^\Sigma_M(G))
 &\overset{\eqref{eq:np-total-depth-added-vertices-bounds},\,
                 \eqref{eq:np-total-depth-W-gap}}{\leq}
 M\bigl((186+6\epsilon)m+3\bigr)+N(N+1)/2,\\
 \textup{NO:}\quad \totd(H^\Sigma_M(G))
 &\overset{\eqref{eq:np-total-depth-added-vertices-bounds},\,
                 \eqref{eq:np-total-depth-W-gap}}{\geq}
 M\bigl((189-6\epsilon)m+3\bigr).
\end{aligned}
\]
The NO lower bound minus $r$ times the YES upper bound is at least
\[
M\bigl(\gamma m-3(r-1)\bigr)-rN(N+1)/2
  \quad\overset{\eqref{eq:np-total-depth-source-scale}}{>}
       M\gamma m/2-rN(N+1)/2
 \quad\overset{\eqref{eq:np-total-depth-amplification-budget}}{\geq}0.
\]
The output has $99m+3M=O(m)$ vertices and $O(m^2)$ edges, and is
constructed deterministically in polynomial time. Its completed core
is connected, because each part is a nonempty clique, auxiliary pairs
join $B$ to $C$, and literal-block edges join $A$ to the other parts.
Every added vertex has a neighbor in its corresponding part, so the
whole output is connected. An $r$-approximation would distinguish the
source YES and NO cases.
\end{proof}
